\chapter{Verification code for the Bosch type theorem}\label{app:bosch}

The script behind the numerical checks of Section~\ref{ch:boschB}, \texttt{code\_bosch/\allowbreak verify\_\allowbreak bosch\_\allowbreak xi\_\allowbreak v7.py}, is in the book's archive (Appendix~\ref{app:code}) with its output, \texttt{verify\_\allowbreak bosch\_\allowbreak xi\_\allowbreak v7\_\allowbreak output.txt}, and the line numbers below are those of these two files; an identical copy of both is in \texttt{bosch\_hcm\_xi/}, beside the authors' notes. The script consists of thirteen blocks that run one after another in one process, each with its own imports; blocks 8 and 9 use \texttt{numpy} as well as \texttt{mpmath}, block 8 passes quadrature data to block 9, and block 13 calls the functions that block 12 defines. Run in an empty folder, the blocks write their results to nine text files, one per block except that block 8 writes none, blocks 5 to 7 share a file and so do blocks 9 and 10, and the output file joins these files in block order. Table~\ref{tab:app:bosch} gives, for each block, its lines in the script and in the output and the statements of Section~\ref{ch:boschB} that it checks.

\begin{table}[htbp]
\centering\footnotesize\setlength{\tabcolsep}{4pt}
\begin{tabular}{@{}lllp{0.37\textwidth}p{0.25\textwidth}@{}}
\toprule
block & script & output & what is checked & statement\\
\midrule
1 & 1 to 50 & 1 to 10 & \raggedright HCM test by the signs of the Taylor coefficients of pair functions, for Bosch's family and for a toy with zeros $14e^{\pm0.1i}$ and $21$; minimum of ${\Im\varepsilon_\varphi(re^{i\theta})/r}$ over $885$ points & \raggedright Theorems~\ref{bs:thm:bosch}, \ref{bs:thm:power}; Corollary~\ref{bs:cor:uncond}\tabularnewline
2 & 51 to 86 & 11 to 24 & \raggedright minimum of $\Im\varepsilon_h$ over $49$ radii and $199$ angles, for Bosch's family and for the toy & \raggedright Theorems~\ref{bs:thm:crit}, \ref{bs:thm:bosch}, \ref{bs:thm:power}\tabularnewline
3 & 87 to 116 & 25 to 36 & \raggedright $u_p$ from \eqref{bs:eq:up} against \eqref{bs:eq:upzeros} with $30$ zeros; $\log\varphi(s^p)$ recovered from $u_p$ & \raggedright Theorem~\ref{bs:thm:ggc}\tabularnewline
none &  & 37 to 39 & \raggedright \eqref{bs:eq:upzeros} with $100$ and $200$ zeros; $u_p>0$ at $148$ points & \raggedright Theorem~\ref{bs:thm:ggc}; Corollary~\ref{bs:cor:rhX}\tabularnewline
4 & 117 to 148 & 40 to 49 & \raggedright density of ${X_C=\sum_ke_k/(2k-1)^2}$ against $(4/\pi)\eta(4it/\pi)^3$, and its Laplace transform at $s=2$; a sufficient inequality in the prime region & \raggedright Theorem~\ref{bs:thm:eta}; Proposition~\ref{bs:prop:prime}\tabularnewline
5, 6 & 149 to 199 & 50 to 60 & \raggedright residue coefficients $c_n$ from $200$ zeros, their signs, sizes and gaps; the residue series against Talbot inversion & \raggedright Theorem~\ref{bs:thm:nb} and the numerical paragraph after it\tabularnewline
7 & 200 to 219 & 61 to 66 & \raggedright the bound $m_p(t)$ of \eqref{bs:eq:primebound}; ${-\zeta'/\zeta(\sigma)<1/(\sigma-1)}$ & \raggedright Proposition~\ref{bs:prop:prime}; Table~\ref{bs:tab:prime}\tabularnewline
8 to 10 & 220 to 318 & 67 to 80 & \raggedright $g$ by quadrature and by series; Pick scan of $-g'/g$; $(\log g)''$ on $[0,1]$ & \raggedright Proposition~\ref{bs:prop:Wcrit}; Theorem~\ref{bs:thm:gap}\tabularnewline
11 & 319 to 343 & 81 to 86 & \raggedright the pair kernel; constants for ${T_0=3\times10^{12}}$ & \raggedright Theorem~\ref{bs:thm:gpos}\tabularnewline
12, 13 & 344 to 419 & 87 to 108 & \raggedright Pick strip $h_*$ and threshold $\beta_*=\pi/h_*$ in logistic models with frequencies removed & \raggedright Table~\ref{bs:tab:gaps}\tabularnewline
none &  & 109 & \raggedright gaps among the first $1000$ zeros & \raggedright Question~\ref{bs:q:gap}\tabularnewline
\bottomrule
\end{tabular}
\caption{The blocks of the verification script: their lines in the script and in the output, what they check, and the statements of Section~\ref{ch:boschB} they support. Output lines 37 to 39 and 109 come from code that is not in the script, and lines 93 to 102 from an earlier version of block 13.}
\label{tab:app:bosch}
\end{table}

The script runs as it stands. In an empty folder it ends after about $400$ seconds, and its nine files, joined in block order, give $95$ lines, each of which occurs in the output file in the same order and agrees with it byte for byte, apart from the timings in seconds that some lines report. The output file is nevertheless not the record of a single run, and its other $14$ lines come from code that is not in the script. Lines 37 to 39 (the zero sums with $100$ and $200$ zeros and the positivity of $u_p$ at $148$ points) and line 109 (the gaps among the first $1000$ zeros) are reproduced exactly by reconstructions kept in the archive (Section~\ref{app:code:repro}), although the grid of $148$ points is a guess. Lines 93 to 102 come from an earlier version of block 13 with ten cases. Lines 103 to 107 run five of these cases again with the present block 13: for the removed frequencies $\{3\}$ they repeat line 102, and for $\{2,3\}$, $\{2,3,5,6\}$, $\{3,4\}$ and $\{4,5\}$ they give $\pi/h_*=2.99997$, $3.0$, $1.48313$ and $1.31671$ where lines 94, 97, 99 and 100 give $1.0$, the value printed when the scan finds no failure below $\pi$. Lines 103 to 106 therefore supersede those four lines; the present block 13, run on all ten cases, reproduces lines 93, 95, 96, 98, 101 and 102, from which Table~\ref{bs:tab:gaps} takes its other rows. Line 108 adds the logistic model with no frequency removed.

Two results in the output need care in reading. For the toy with zeros $14e^{\pm0.1i}$ and $21$, neither test locates the threshold $p_*=0.9363$. The sign test finds no violation at $p=0.95$ (line 8), and the scan of block 2 misses the poles at $p=0.94$ and $p=0.97$, which lie between its grid points or beyond its last ray (lines 22 and 23, where the flag for poles is computed from their position, not detected); both tests see the failure only at $p=1$ (lines 9 and 24). At $t=0.003$ the terms of the residue series reach $2.7\times10^{20}$ while the density is about $3.3\times10^{-72}$, so about $92$ digits cancel, more than the $25$ carried in line 56 and the $90$ carried in line 60: the series values at $t=0.003$ in these lines are rounding error, and the Talbot values there are negative. At $150$ digits the residue series over the first $250$ zeros and Talbot inversion agree on the density $3.32007507753\times10^{-72}$ at $t=0.003$, with a largest term of $2.694\times10^{20}$ and $91.9$ digits cancelled; the reconstructions of Section~\ref{app:code:repro} contain this computation.
